% GENERATED FILE. Edit canonical lessons, then run npm run generate:books.
% canonical-source-hash: fnv1a64:c1bc22491aa4bc5b
% canonical-lessons: GE-C02-ich, GE-C02-du-sie, GE-C02-heissen, GE-C02-heissen-endungen, GE-C02-ich-heisse, GE-C02-wie, GE-C02-wie-heissen-sie, GE-C02-mich, GE-C02-freut-mich, GE-C02-practice

\chapter{Introducing Yourself}
\label{ch:introductions}
\hlchaptermodality{2}

\begin{chapteropening}
\textbf{By the end of this chapter:} \emph{I can introduce myself in German, give my name, and understand the other person's name.}

Now a first real exchange — giving your name, asking for someone else's, and saying you are glad to meet them: Ich heiße Mira. — Wie heißen Sie? — Ich heiße Arun. — Freut mich. It is built the way the greetings were: each piece taken apart first, then assembled — and, because German is English's sibling, most of those pieces are direct English cousins. GE-C02-practice puts the four lines back together at the end.
\end{chapteropening}

\section[ich]{ich — \textquotedblleft{}I,\textquotedblright{} the same ancient word as Latin \emph{ego}}

\label{lesson:GE-C02-ich}



\begin{quote}
To introduce yourself, first \textquotedblleft{}I.\textquotedblright{} German's is a three-letter word hiding a very deep root.
\end{quote}

\begin{sounds}
\begin{itemize}
\raggedright
  \item \textbf{ich} = \emph{ikh} — the \emph{ch} here is \textbf{soft}: a light hiss at the front of the mouth (the \textquotedblleft{}ich-sound\textquotedblright{}), \emph{not} the throat-scrape of \emph{Nacht}. German has both; after \emph{i} and \emph{e} it's the soft one.
\end{itemize}
\end{sounds}

\begin{cousinweb}
\textbf{ich} — \textquotedblleft{}I.\textquotedblright{} Root: Proto-Germanic \textbf{*ik}, from Proto-Indo-European \textbf{*eǵ} — the \emph{same} deep root that became:

\begin{itemize}
\raggedright
  \item Latin \textbf{ego} (and English \textbf{ego}), Greek \emph{egō},
  \item and, worn down hardest of all, English \textbf{I}.
\end{itemize}

So German \emph{ich}, Latin \emph{ego}, and English \emph{I} are three survivals of \textbf{one} ancient word for \textquotedblleft{}I\textquotedblright{} — English filed it down to a single letter, German kept a bit more.
\end{cousinweb}

\subsection*{Your turn}
Say these aloud:

\begin{itemize}
\raggedright
  \item \textquotedblleft{}ich\textquotedblright{} (\emph{ikh}, soft \emph{ch})
  \item the three cousins — \emph{ich} / \emph{ego} / \emph{I}
\end{itemize}

\begin{tcolorbox}[breakable,colback=teal!4,colframe=teal!35!black,title={Before you move on}]
What PIE root gave \textbf{ich}, and what Latin and English words are its cousins? (\emph{*eǵ} — Latin \emph{ego}, English \emph{I}.) Which \emph{ch} sound does \emph{ich} use — soft or the \emph{Nacht} scrape? (The soft \textquotedblleft{}ich-sound.\textquotedblright{})
\end{tcolorbox}
\section[du / Sie]{du / Sie — \textquotedblleft{}you,\textquotedblright{} familiar and formal}

\label{lesson:GE-C02-du-sie}



\begin{quote}
To ask someone's name you need \textquotedblleft{}you\textquotedblright{} — and German's formal one is the strangest of all the languages here.
\end{quote}

\begin{sounds}
\begin{itemize}
\raggedright
  \item \textbf{du} = \emph{doo} (long \emph{u}, as in \emph{boot}). \textbf{Sie} = \emph{zee} (the \emph{s} is a \emph{z}-sound at the start).
\end{itemize}
\end{sounds}

\begin{cousinweb}
\textbf{du} (familiar \textquotedblleft{}you\textquotedblright{}) $\leftarrow$ Proto-Germanic \textbf{*þū}, from PIE \textbf{*tū} — the \emph{same} root as:

\begin{itemize}
\raggedright
  \item English archaic \textbf{thou}, Latin \textbf{tū}, French \textbf{tu}.
\end{itemize}

English kept \emph{thou} only in old prayers and some dialects; German uses \emph{du} every day.
\end{cousinweb}

\begin{grammarlens}[title={Grammar Lens: two \textquotedblleft{}you\textquotedblright{}s — and German's is the odd one out}]
\begin{itemize}
\raggedright
  \item \textbf{du} — friends, family, children, peers. Intimate, equal.
  \item \textbf{Sie} — the respectful \textquotedblleft{}you,\textquotedblright{} for strangers and elders.
\end{itemize}

Here's the twist: \textbf{Sie} is really the word for \textbf{\textquotedblleft{}they\textquotedblright{}} (\emph{sie}), promoted to a polite \textquotedblleft{}you\textquotedblright{} and always \textbf{Capitalized} (\emph{Sie}) to tell it apart from \emph{sie} = \textquotedblleft{}they/she.\textquotedblright{} So German is polite by addressing you in the \textbf{third person plural} — as if saying \textquotedblleft{}might \emph{they}…\textquotedblright{} about you to your face. Three languages, three escapes from the blunt \textquotedblleft{}you\textquotedblright{}:

\begin{itemize}
\raggedright
  \item Spanish coined \textbf{usted} (\textquotedblleft{}your grace\textquotedblright{}),
  \item French uses the \textbf{2nd-person plural} \emph{vous} (\textquotedblleft{}you all\textquotedblright{}),
  \item German borrows the \textbf{3rd-person plural} \emph{Sie} (\textquotedblleft{}they\textquotedblright{}).
\end{itemize}

When unsure, use \textbf{Sie}.
\end{grammarlens}

\subsection*{Your turn}
Say these aloud:

\begin{itemize}
\raggedright
  \item \textquotedblleft{}du\textquotedblright{} (\emph{doo}), \textquotedblleft{}Sie\textquotedblright{} (\emph{zee})
  \item who gets \emph{du} (friends), who gets \emph{Sie} (strangers, elders)
  \item what \emph{Sie} literally is (capitalized \textquotedblleft{}they,\textquotedblright{} used as polite \textquotedblleft{}you\textquotedblright{})
\end{itemize}

\begin{tcolorbox}[breakable,colback=teal!4,colframe=teal!35!black,title={Before you move on}]
What is the formal \textbf{Sie} literally, and why capitalized? (It's \textquotedblleft{}they,\textquotedblright{} \emph{sie}, promoted to polite \textquotedblleft{}you\textquotedblright{}; capitalized to distinguish it.) What old English word was the lost \emph{du}? (\emph{thou}.)
\end{tcolorbox}
\section[heißen]{heißen — \textquotedblleft{}to be called,\textquotedblright{} the verb English almost lost}

\label{lesson:GE-C02-heissen}



\begin{quote}
German gives its name with one plain verb — \emph{heißen}, \textquotedblleft{}to be called.\textquotedblright{} No \textquotedblleft{}myself\textquotedblright{} needed, unlike the Romance languages.
\end{quote}

\begin{sounds}
\begin{itemize}
\raggedright
  \item \textbf{heißen} = \emph{HY-ssen}. The \textbf{ß} (\emph{Eszett}) is a sharp \textquotedblleft{}ss.\textquotedblright{} \textbf{ei} is the vowel of English \emph{eye}. So \textbf{heiße} = \emph{HY-ssuh}.
\end{itemize}
\end{sounds}

\begin{cousinweb}
\textbf{heißen} — \textquotedblleft{}to be called / to be named.\textquotedblright{} Root: Proto-Germanic \textbf{*haitaną}. English \emph{had} this verb:

\begin{itemize}
\raggedright
  \item \textbf{hight}, \textquotedblleft{}named\textquotedblright{} — now archaic (\textquotedblleft{}a knight \emph{hight} Gawain\textquotedblright{}), and
  \item \textbf{behest}, \textquotedblleft{}a command\textquotedblright{} — literally something \emph{called for}.
\end{itemize}

So German kept as an everyday word a verb English let fade to poetry. And where Spanish and French say \textquotedblleft{}I call \emph{myself}\textquotedblright{} (a reflexive), German uses \emph{heißen} as a plain verb of \textbf{being} called — no \textquotedblleft{}myself\textquotedblright{} at all.
\end{cousinweb}

\subsection*{Your turn}
Say these aloud:

\begin{itemize}
\raggedright
  \item \textquotedblleft{}heißen\textquotedblright{} (\emph{HY-ssen})
  \item the English fossil — \emph{hight} (\textquotedblleft{}named\textquotedblright{}), \emph{behest}
\end{itemize}

\begin{tcolorbox}[breakable,colback=teal!4,colframe=teal!35!black,title={Before you move on}]
What does \textbf{heißen} mean, and what archaic English word is its cousin? (\textquotedblleft{}To be called\textquotedblright{}; \emph{hight}.) How is German's way of giving a name different from Spanish/French? (A plain verb \textquotedblleft{}to be called,\textquotedblright{} no reflexive \textquotedblleft{}myself.\textquotedblright{})
\end{tcolorbox}
\section[ich heiße / du heißt / Sie heißen]{heißen with a person on the front — one verb, three endings}

\label{lesson:GE-C02-heissen-endungen}



\begin{quote}
You have the verb and you have the people. \emph{heißen} is the form a dictionary prints; it is not the form you say. Put a person in front of it and the tail of the word moves.
\end{quote}

\begin{grammarlens}[title={Grammar Lens: one verb, changing its ending}]
\textbf{heißen} is the dictionary (infinitive) form. German swaps the ending for each person:

\begin{itemize}
\raggedright
  \item \textbf{ich heiße} — I am called
  \item \textbf{du heißt} — you (familiar) are called
  \item \textbf{Sie heißen} — you (formal) are called
\end{itemize}

That ending-swap is the same machine on every German verb. You'll use \emph{ich heiße} and \emph{Sie heißen} in the next lessons.
\end{grammarlens}

\subsection*{Your turn}
Say these aloud:

\begin{itemize}
\raggedright
  \item the three forms — ich heiße / du heißt / Sie heißen
  \item the infinitive, then the \emph{ich} form — heißen, ich heiße
  \item the tails alone — \emph{-e}, \emph{-t}, \emph{-en}
\end{itemize}

\begin{tcolorbox}[breakable,colback=teal!4,colframe=teal!35!black,title={Before you move on}]
Which ending goes with \emph{ich}? (\emph{-e} — \emph{ich heiße}.) With formal \emph{Sie}? (\emph{-en} — \emph{Sie heißen}, the same as the dictionary form.) Is this ending machine special to \emph{heißen}? (No — every German verb runs it.)
\end{tcolorbox}
\section[ich heiße…]{ich heiße… — \textquotedblleft{}my name is…,\textquotedblright{} literally \textquotedblleft{}I am called\textquotedblright{}}

\label{lesson:GE-C02-ich-heisse}



\begin{quote}
Assemble the two pieces — \emph{ich} and \emph{heißen} — into the sentence you need for meeting anyone.
\end{quote}

\subsection*{You'll want to know: ich heiße…}
\begin{itemize}
\raggedright
  \item \textbf{ich} (I) + \textbf{heiße} (am called).
\end{itemize}

\begin{quote}
\textbf{ich heiße…} = word for word \textquotedblleft{}I am called…\textquotedblright{} = \textbf{\textquotedblleft{}my name is…\textquotedblright{}}
\end{quote}

\emph{Ich heiße Arun.} $\to$ \textquotedblleft{}My name is Arun.\textquotedblright{} \emph{Ich heiße Mira.} $\to$ \textquotedblleft{}My name is Mira.\textquotedblright{}

\begin{culture}
German also has the literal \textquotedblleft{}my name is\textquotedblright{}:

\begin{itemize}
\raggedright
  \item \textbf{mein Name ist…} — where \textbf{Name} $\leftarrow$ Proto-Germanic \textbf{*namô}, the \emph{same} root as English \textbf{name} and Latin \textbf{nōmen} ($\to$ English \textbf{noun}, \textbf{nom}inate). It's correct and common.
\end{itemize}

But \textbf{ich heiße} is warmer and more idiomatic — German prefers the \textbf{verb} (\textquotedblleft{}I am called\textquotedblright{}) over the \textbf{noun} phrase (\textquotedblleft{}my name is\textquotedblright{}). (Note \textbf{Name} is Capitalized: German capitalizes \emph{every} noun — you'll see it on every naming word.)
\end{culture}

\subsection*{Your turn}
Say these aloud:

\begin{itemize}
\raggedright
  \item \textquotedblleft{}ich heiße…\textquotedblright{} then your own name
  \item what it literally says (\textquotedblleft{}I am called\textquotedblright{})
  \item the literal alternative — \textquotedblleft{}mein Name ist…\textquotedblright{}
\end{itemize}

\begin{tcolorbox}[breakable,colback=teal!4,colframe=teal!35!black,title={Before you move on}]
What does \textbf{ich heiße} literally mean? (\textquotedblleft{}I am called.\textquotedblright{}) What's the more literal \textquotedblleft{}my name is,\textquotedblright{} and what English words is \emph{Name} cousin to? (\emph{mein Name ist}; \emph{name}, \emph{noun}.)
\end{tcolorbox}
\section[wie]{wie — \textquotedblleft{}how,\textquotedblright{} cousin of English \emph{how}}

\label{lesson:GE-C02-wie}



\begin{quote}
German asks \textquotedblleft{}what's your name?\textquotedblright{} as \textquotedblleft{}\textbf{how} are you called?\textquotedblright{} — so first, \textquotedblleft{}how.\textquotedblright{}
\end{quote}

\begin{sounds}
\begin{itemize}
\raggedright
  \item \textbf{wie} = \emph{vee}. Two traps: German \textbf{w} is an English \emph{v}, and \textbf{ie} is a long \emph{ee}. So \emph{wie} sounds like English \textquotedblleft{}vee,\textquotedblright{} never \textquotedblleft{}wy.\textquotedblright{}
\end{itemize}
\end{sounds}

\begin{cousinweb}
\textbf{wie} — \textquotedblleft{}how.\textquotedblright{} Root: Proto-Germanic \textbf{*hwī}, from the PIE interrogative root \textbf{*k\textsuperscript{w}o-} — the very root behind the English \textbf{wh-} words:

\begin{itemize}
\raggedright
  \item \textbf{how}, \textbf{what}, \textbf{who}, \textbf{why} — and, through Latin, the \textbf{qu-} of \emph{question}, \emph{quantity}, and the \emph{quo modo} that became French \emph{comment}.
\end{itemize}

German \emph{wie} and English \emph{how} are the same question-word, lightly reshaped.
\end{cousinweb}

\begin{grammarlens}[title={Grammar Lens: German counts names as a matter of \textquotedblleft{}how\textquotedblright{}}]
Like French, German asks a name with \textquotedblleft{}\textbf{how},\textquotedblright{} not \textquotedblleft{}what\textquotedblright{}: \emph{wie heißen Sie?} is literally \textquotedblleft{}\emph{how} are you called?\textquotedblright{} — about the \emph{manner} of your being called. Don't translate word for word; learn the logic: the name-question runs on \textbf{wie}.
\end{grammarlens}

\subsection*{Your turn}
Say these aloud:

\begin{itemize}
\raggedright
  \item \textquotedblleft{}wie\textquotedblright{} (\emph{vee} — \emph{w} is \emph{v}, \emph{ie} is \emph{ee})
  \item the English \emph{wh-} cousins — how, what, who, why
  \item \textquotedblleft{}how are you called?\textquotedblright{} is the German for \textquotedblleft{}what's your name?\textquotedblright{}
\end{itemize}

\begin{tcolorbox}[breakable,colback=teal!4,colframe=teal!35!black,title={Before you move on}]
How is \textbf{wie} pronounced, and what English question-words share its root? (\emph{vee}; how / what / who / why.) English asks \textquotedblleft{}what is your name?\textquotedblright{} — what does German literally ask? (\textquotedblleft{}\emph{How} are you called?\textquotedblright{})
\end{tcolorbox}
\section[wie heißen Sie?]{wie heißen Sie? — \textquotedblleft{}what's your name?\textquotedblright{}}

\label{lesson:GE-C02-wie-heissen-sie}



\begin{quote}
Assemble it: \textquotedblleft{}how\textquotedblright{} + \textquotedblleft{}are you called\textquotedblright{} = the polite question that opens the introduction.
\end{quote}

\subsection*{You'll want to know: wie heißen Sie?}
\begin{itemize}
\raggedright
  \item \textbf{wie} (how) + \textbf{heißen Sie} (\textquotedblleft{}are you {[}formal{]} called\textquotedblright{}).
\end{itemize}

\begin{quote}
\textbf{wie heißen Sie?} = literally \textquotedblleft{}\textbf{how are you called?}\textquotedblright{} = the polite \textbf{\textquotedblleft{}what's your name?\textquotedblright{}}
\end{quote}

\begin{grammarlens}[title={Grammar Lens: verb in second place}]
German likes the verb in the \textbf{second} slot of a sentence:

\begin{quote}
\emph{Wie} (1) \emph{heißen} (2) \emph{Sie}?
\end{quote}

The \textbf{informal} version swaps in \emph{du} and its ending:

\begin{quote}
\emph{Wie \textbf{heißt du}?}
\end{quote}

So the one verb \emph{heißen} runs the whole exchange: \emph{ich heiße…} to give your name, \emph{wie heißen Sie?} to ask for someone else's. (Reply to \emph{wie heißen Sie?} with \emph{Ich heiße…}.)
\end{grammarlens}

\subsection*{Your turn}
Say these aloud:

\begin{itemize}
\raggedright
  \item \textquotedblleft{}wie heißen Sie?\textquotedblright{} (formal)
  \item the informal — \textquotedblleft{}wie heißt du?\textquotedblright{}
  \item the pair — ask, then answer: \textquotedblleft{}wie heißen Sie?\textquotedblright{} / \textquotedblleft{}ich heiße…\textquotedblright{}
\end{itemize}

\begin{tcolorbox}[breakable,colback=teal!4,colframe=teal!35!black,title={Before you move on}]
What does \textbf{wie heißen Sie?} literally say? (\textquotedblleft{}How are you called?\textquotedblright{}) What's the informal version, and where does German put the verb? (\emph{Wie heißt du?}; the verb goes second.)
\end{tcolorbox}
\section[mich]{mich — \textquotedblleft{}me,\textquotedblright{} the same word as \emph{ich}, doing a different job}

\label{lesson:GE-C02-mich}



\begin{quote}
English does not say \textquotedblleft{}gladdens I.\textquotedblright{} It says \textquotedblleft{}gladdens \textbf{me}.\textquotedblright{} German makes the same move, and the next phrase you learn needs it.
\end{quote}

\subsection*{You'll want to know: mich}
\begin{quote}
\textbf{mich} — \textquotedblleft{}me,\textquotedblright{} the object form of \emph{ich}.
\end{quote}

Said \emph{mikh}: the final \textbf{ch} is the soft ich-sound you already met in \emph{ich}, not the throat-scrape of \emph{Nacht}.

\begin{cousinweb}
\textbf{mich} (\textquotedblleft{}me,\textquotedblright{} the object form) $\leftarrow$ Proto-Germanic \textbf{*mek}, from PIE \textbf{*me-} — the \emph{same} first-person root as English \textbf{me}, \textbf{my}, \textbf{mine} (and the French \emph{me} in \emph{je m'appelle}). So German \emph{mich}, English \emph{me}, and French \emph{me} are three cousins of one ancient word for \textquotedblleft{}me.\textquotedblright{} And just as English shifts \textbf{I} (subject) $\to$ \textbf{me} (object), German shifts \textbf{ich} $\to$ \textbf{mich}.
\end{cousinweb}

\subsection*{Your turn}
Say these aloud:

\begin{itemize}
\raggedright
  \item \textquotedblleft{}ich\textquotedblright{} then \textquotedblleft{}mich\textquotedblright{} — subject, then object
  \item the three cousins — \emph{mich} / \emph{me} / \emph{me}
  \item \textquotedblleft{}mich\textquotedblright{} (\emph{mikh}, soft \emph{ch} — not the \emph{Nacht} scrape)
\end{itemize}

\begin{tcolorbox}[breakable,colback=teal!4,colframe=teal!35!black,title={Before you move on}]
What is \emph{mich}, and what is its subject form? (\textquotedblleft{}Me\textquotedblright{}; \emph{ich}.) Which English pair does \emph{ich $\to$ mich} match? (\emph{I $\to$ me}.) Which \emph{ch} does \emph{mich} use? (The soft ich-sound.)
\end{tcolorbox}
\section[freut mich]{freut mich — \textquotedblleft{}pleased to meet you,\textquotedblright{} i.e. \textquotedblleft{}it gladdens me\textquotedblright{}}

\label{lesson:GE-C02-freut-mich}



\begin{quote}
The warm close to an introduction — and it literally says the meeting makes you glad.
\end{quote}

\begin{sounds}
\begin{itemize}
\raggedright
  \item \textbf{freut mich} = \emph{froyt mikh}. \textbf{eu} is the vowel of English \emph{boy}; the final \textbf{ch} is the soft ich-sound again (from \emph{ich}).
\end{itemize}
\end{sounds}

\begin{cousinweb}
\textbf{freut mich} is short for \textbf{es freut mich (Sie kennenzulernen)} — \textquotedblleft{}it \textbf{gladdens} me (to get to know you).\textquotedblright{} The verb \textbf{freuen} (\textquotedblleft{}to gladden\textquotedblright{}) sits on \textbf{froh} (\textquotedblleft{}glad, joyful\textquotedblright{}):

\begin{itemize}
\raggedright
  \item the same \textbf{froh} in \textbf{Frohe Weihnachten} (\textquotedblleft{}Merry Christmas\textquotedblright{}), and a cousin of English \textbf{frolic} (via Dutch \emph{vrolijk}, \textquotedblleft{}merry\textquotedblright{}).
\end{itemize}

So \emph{freut mich} says, warmly, \textquotedblleft{}this makes me glad.\textquotedblright{}
\end{cousinweb}

\begin{grammarlens}[title={Grammar Lens: \textquotedblleft{}it gladdens \emph{me}\textquotedblright{}}]
German literally says \textquotedblleft{}{[}it{]} gladdens \textbf{me}\textquotedblright{} — the \emph{meeting} is what does the gladdening, and you (\textbf{mich}) are on the receiving end. Two common variants: \textbf{sehr erfreut} (\textquotedblleft{}much pleased\textquotedblright{}) and the plain \textbf{angenehm} (\textquotedblleft{}pleasant / agreeable\textquotedblright{}), both usable as \textquotedblleft{}pleased to meet you.\textquotedblright{}
\end{grammarlens}

\subsection*{Your turn}
Say these aloud:

\begin{itemize}
\raggedright
  \item \textquotedblleft{}freut mich\textquotedblright{} (\emph{froyt mikh})
  \item the root — \emph{froh} (\textquotedblleft{}glad\textquotedblright{}), as in \emph{Frohe Weihnachten}
  \item a variant — \textquotedblleft{}sehr erfreut\textquotedblright{} or \textquotedblleft{}angenehm\textquotedblright{}
\end{itemize}

\begin{tcolorbox}[breakable,colback=teal!4,colframe=teal!35!black,title={Before you move on}]
What does \textbf{freut mich} literally mean, and on what word is \emph{freuen} built? (\textquotedblleft{}It gladdens me\textquotedblright{}; \emph{froh}, \textquotedblleft{}glad.\textquotedblright{}) What does \emph{mich} mean in it? (\textquotedblleft{}Me\textquotedblright{} — the meeting gladdens \emph{me}.)
\end{tcolorbox}
\section[Practice]{Chapter 2 — the introduction, whole}

\label{lesson:GE-C02-practice}



\begin{quote}
Every piece is built. Here's the exchange it was all for.
\end{quote}

\subsection*{The exchange}
\noindent
\begin{tabularx}{\linewidth}{@{}>{\raggedright\arraybackslash}X>{\raggedright\arraybackslash}X@{}}
\toprule
\textbf{German} & \textbf{English} \\
\midrule
\emph{Ich heiße Mira.} & My name is Mira. \\
\emph{Wie heißen Sie?} & What's your name? \\
\emph{Ich heiße Arun.} & My name is Arun. \\
\emph{Freut mich.} & Pleased to meet you. \\
\bottomrule
\end{tabularx}

Every piece a Germanic cousin of English:

\begin{itemize}
\raggedright
  \item \textbf{ich} $\leftarrow$ \emph{*ik} (English \emph{I}, Latin \emph{ego})
  \item \textbf{mich} $\leftarrow$ \emph{*mek} (English \emph{me}, \emph{my}, \emph{mine})
  \item \textbf{heißen} $\leftarrow$ \emph{*haitaną} (English \emph{hight}, \emph{behest})
  \item \textbf{wie} $\leftarrow$ \emph{*hwī} (English \emph{how}, \emph{what}, \emph{who})
  \item \textbf{freuen} $\leftarrow$ \emph{froh} (glad; English \emph{frolic})
\end{itemize}

And two habits of German to keep: it names with a \textbf{plain verb of being called} (\textquotedblleft{}I am called,\textquotedblright{} no \textquotedblleft{}myself\textquotedblright{}), and it is polite by the \textbf{third person} (\emph{Sie} = \textquotedblleft{}they\textquotedblright{}).

\subsection*{Your turn}
Say these aloud:

\begin{itemize}
\raggedright
  \item the whole exchange, both voices
  \item your own introduction — \textquotedblleft{}ich heiße …, freut mich\textquotedblright{}
  \item ask it back — \textquotedblleft{}wie heißen Sie?\textquotedblright{}
\end{itemize}

\begin{tcolorbox}[breakable,colback=teal!4,colframe=teal!35!black,title={Before you move on}]
Give your name in German, ask someone else's, and say you're pleased to meet them. (\emph{Ich heiße… / Wie heißen Sie? / Freut mich.}) What two German habits does this show? (Naming with a plain \textquotedblleft{}to be called\textquotedblright{} verb; politeness by the third person, \emph{Sie}.)

Next chapter: \emph{Wie geht's?} (\textquotedblleft{}how's it going?\textquotedblright{}) and answering with \emph{gut} again — the responding cycle.
\end{tcolorbox}
